\section*{序言（食用指南）}


本讲义主要参考北大版《高等数学》（第三版），
按照高等数学B课程要求编写，
可以配合教材作为学习和复习的辅助材料使用.
讲义分为十二章，
与教材的章节划分一致.
每一章主要分为“\textbf{复习要点}”和“\textbf{解题方法}”两大部分.
\vspace{1em}

“\textbf{复习要点}”主要安排以下内容：

(1)~~构建知识框架；

(2)~~梳理需要集中背诵的定理、公式等；

(3)~~深入展开本章需要加强理解的知识点、辨析易混淆的概念、提醒易错点；

(4)~~拓展相关知识（小节标题以“$*$”标注，若不感兴趣可自由跳过）.\vspace{1em}

“\textbf{解题方法}”主要安排以下内容：

(1)~~重点题型的解题方法梳理
（题型是否“重点”是以北京大学高数B期中/期末考试为标准的）；

(2)~~例题和解答过程，有时附有分析，展示思考过程；

(3)~~进阶的解题方法和例题（小节标题以“*”标注）\vspace{1em}

这份讲义的编写和使用历经了一年的时间.
在此期间，
我的数学水平、
表达能力和排版技术还处于进步的过程中，
也不断发现此前已经落笔的章节还有许多错漏和缺憾.
在读者们的帮助下，
我会在后续的版本中继续修正错误.
也诚恳地希望读者们对讲义提出宝贵意见，
以便改进.

这份讲义能够编写完成，
离不开许多人的帮助和共同努力.
感谢北京大学生命科学学院学业促进中心提供的机会和支持.
感谢与我合作的另两位高等数学学业辅导员卢卓然和王子烨，
他们准备的讲义和讲课方法也为我提供了重要的参考.
最后，还要感谢2024秋和2025春两届高等数学学业辅导班的全体同学.
在讲义尚未完全成型的时候，
他们充当了“小白鼠”，
不仅向我反馈讲义的使用体验，
而且还帮助我改正讲义中的错误，
提供了许多有用的建议.
这次高数辅导的契机让我时隔两年又“重学”了一遍高数.
翻阅课本的时候，
我经常想起自己在大一学习高数的情景，
想起自己也曾对着题目纠结万分，
向同学寻求过帮助，
有时还演变成激烈的争论.
我原本只会成为这些情景的回忆者和旁观者.
但是由于高数辅导，由于编写这份讲义，
我也幸运地成为了记录者.
我相信，
同学们在学术问题上的纯粹，
以及相互包容、信任、合作、携手进步的精神，
都在增加大学生活和学术道路上的适合度.
这也是我所记录下来的，
北大在第一个学期给同学们带来的重要一课.


\begin{flushright}
    2025.10.16
\end{flushright}